\section{Related Work}
\label{sec:related_work}

\subsection{Flow Matching for Robot Policies}

Flow matching~\cite{lipman2023flow,liu2023flow} learns a continuous-time velocity field $v_\theta(x_t, t)$ that transports samples from a prior distribution to the data distribution via an ODE $\mathrm{d}x_t = v_\theta(x_t, t)\,\mathrm{d}t$. Compared to diffusion models~\cite{ho2020denoising,song2021score}, flow matching offers simpler training (no noise schedule), deterministic sampling, and straighter transport paths when combined with optimal transport (OT-CFM)~\cite{tong2024improving}.

In robotics, Diffusion Policy~\cite{chi2023diffusion} first demonstrated that denoising diffusion can effectively model multi-modal action distributions for visuomotor control, introducing the action chunking paradigm where a policy predicts a sequence of future actions rather than a single action. VITA~\cite{gao2026vita} extended this idea to flow matching, proposing to flow directly from visual observation latents to action latents---eliminating the random noise starting point and enabling natural multi-modal action generation. $\pi_0$~\cite{black2024pi0} scaled flow matching policies to foundation-model-level with vision-language-action pretraining, demonstrating strong zero-shot generalization across tasks. Despite their generative power, all these methods treat each timestep independently, lacking mechanisms for temporal reasoning across the episode. Our work builds upon the VITA architecture and augments it with temporal awareness and self-supervised spatial attention.

\subsection{Temporal Encoding in Imitation Learning}

Incorporating temporal information into imitation learning policies has been explored through several complementary approaches.

\paragraph{Observation stacking and frame differencing.} The simplest strategy is \emph{observation stacking}: concatenating the most recent $k$ frames as input~\cite{mnih2015human}. While straightforward, this provides only local temporal context, does not explicitly encode \emph{what changed} between frames, and scales poorly with horizon length. Our temporal feature differencing (Section~\ref{sec:temporal_diff}) addresses this by computing explicit feature-space differences, encoding ``semantic velocity'' without increasing dimensionality with observation horizon.

\paragraph{Transformer-based temporal modeling.} Action Chunking with Transformers (ACT)~\cite{zhao2023act} uses a transformer encoder-decoder architecture over observation sequences, implicitly capturing temporal dependencies through self-attention. RT-2~\cite{brohan2023rt2} leverages large vision-language models with autoregressive action generation, inheriting the transformer's ability to attend over long contexts. However, these approaches introduce substantial computational overhead and do not explicitly model task progress or phase transitions---temporal structure must be learned implicitly from the attention patterns.

\paragraph{Explicit temporal structure.} Temporal Contrastive Networks (TCN)~\cite{sermanet2018time} learn time-aware embeddings through contrastive learning across video frames from multiple viewpoints, requiring multi-view correspondence data. Hierarchical Diffusion Policy~\cite{ma2024hierarchical} decomposes tasks into sub-goals at different temporal scales, but requires explicit sub-goal definitions. Latent Plans~\cite{lee2024behavior} learn to generate latent action plans with phase-conditioned policies, but rely on pre-defined temporal structures or discrete phase labels.

Our approach differs fundamentally: we extract temporal information \emph{implicitly} from the observation stream through self-predictive dynamics and ordinal ranking, without any manual temporal labels, phase annotations, or multi-view data. The prediction error of our latent dynamics model provides a free, self-supervised phase-transition signal, while the ordinal ranking loss requires only the weak assumption that frames within an episode are temporally ordered.

\subsection{Spatial Attention in Visuomotor Policies}

Spatial attention mechanisms aim to guide the visual encoder to focus on task-relevant regions rather than treating all spatial locations equally.

\paragraph{Learned spatial representations.} Spatial softmax~\cite{levine2016end} computes per-channel expected coordinates from convolutional feature maps, providing a compact spatial representation that has become a standard tool in visuomotor policy learning. However, it applies uniform attention across all channels and does not distinguish between task-relevant and irrelevant regions.

\paragraph{Channel and spatial attention modules.} Squeeze-and-Excitation (SE) networks~\cite{hu2018squeeze} learn channel attention weights through global pooling and lightweight MLPs, re-weighting feature channels by their importance. CBAM~\cite{woo2018cbam} extends this with an additional spatial attention branch that learns a 2D attention map. Coordinate Attention~\cite{hou2021coordinate} factorizes channel attention into two spatial directions, capturing long-range positional dependencies along each axis. While effective for general vision tasks, these modules learn \emph{static} attention patterns from the data distribution; they do not adapt to the current interaction dynamics within an episode.

\paragraph{Knowledge distillation from foundation models.} Vision foundation models such as Grounding DINO~\cite{liu2023grounding} and CLIP~\cite{radford2021learning} encode rich spatial and semantic knowledge that can be transferred to lightweight backbones via attention heatmap matching or feature distillation losses. While effective, such approaches introduce dependencies on specific teacher models, require offline feature extraction pipelines, and produce static spatial targets that do not change with the task phase.

Our \fasa{} module takes a fundamentally different approach: rather than relying on external spatial supervision or learning static attention patterns, it discovers task-relevant regions \emph{self-supervisedly} through inter-frame feature temporal contrast. The key insight is that regions undergoing the most visual change between consecutive frames correspond to active interaction zones---exactly the regions most relevant for manipulation tasks. This makes \fasa{}'s attention inherently \emph{dynamic}, adapting in real time as the task progresses from approach to tracking.

\subsection{Memory-Augmented Policies}

External memory mechanisms have a long history in sequence modeling. Neural Turing Machines~\cite{graves2014neural} and Differentiable Neural Computers~\cite{graves2016hybrid} provide general-purpose read-write memory with content-based and location-based addressing, but their complexity limits practical adoption in real-time robot control. End-to-End Memory Networks~\cite{sukhbaatar2015end} simplify the architecture with multi-hop attention over memory slots, demonstrating effectiveness on question answering and dialogue tasks.

In robot learning, transformer-based policies~\cite{zhao2023act,brohan2023rt2} implicitly maintain memory through the KV-cache mechanism during autoregressive generation, but this couples memory with the specific transformer architecture. Test-Time Training (TTT) layers~\cite{sun2024ttt} replace attention with a learned update rule that compresses context into model weights, achieving linear complexity---but introduce significant additional training complexity and are not straightforward to integrate with MLP-based flow networks.

Our Episodic Memory Bank is deliberately lightweight: a single key-value attention read operation that runs in $O(T)$ time with respect to episode length $T$. It is architecture-agnostic (compatible with any base model that produces observation latents), adds minimal parameters, and requires no modifications to the flow network itself. The key design choice of using same-episode pseudo-memory during training (rather than random batch sampling) reduces the distribution gap between training and inference, where the memory buffer contains genuine trajectory history.

\subsection{Self-Predictive Learning}

Self-predictive representation learning~\cite{grill2020bootstrap,bardes2022vicreg,chen2021exploring} trains encoders by predicting representations of augmented or temporally adjacent views without negative samples. BYOL~\cite{grill2020bootstrap} uses a momentum encoder and stop-gradient target to prevent representation collapse, while SimSiam~\cite{chen2021exploring} demonstrates that a simple stop-gradient predictor suffices without momentum. VICReg~\cite{bardes2022vicreg} takes an alternative approach, regularizing the variance, invariance, and covariance of embeddings.

In reinforcement learning, forward dynamics models have been used both for learning task-relevant representations and for driving exploration. Curiosity-driven exploration~\cite{pathak2017curiosity} uses the prediction error of a learned forward model as an intrinsic reward signal, encouraging the agent to visit novel states. Dreamer~\cite{hafner2020dream} learns a latent world model for planning by imagination, using the world model's predictions to generate synthetic training trajectories.

We draw on this line of work in two novel ways. First, rather than using self-prediction solely for representation learning or exploration, we use the \emph{prediction error} of a latent dynamics model as a feature for downstream progress estimation. The key insight is that prediction error naturally peaks at phase transitions---moments when the observation distribution shifts---providing a free, self-supervised signal for temporal structure discovery. Our dynamics model uses a stop-gradient target projection (inspired by BYOL) to prevent collapse while maintaining meaningful prediction errors. Second, our \fasa{} module extends the self-predictive principle to the spatial domain: instead of predicting future features, it detects \emph{where} features change between frames, using inter-frame temporal contrast as a self-supervised spatial attention signal.
